% Standalone section for inclusion in the main report.
% Requires amsmath, amssymb, and bibliography keys Salehi2002, Wehmeier1997.
\section{What constructive finiteness would assert}
\label{sec:logical-audit}

The word ``finite'' requires particular care in this problem.  Several
classically equivalent statements about integral points have different
intuitionistic meanings.  A proof of one of them must not silently be advertised
as a proof of another.  In particular, a bound on the \emph{number} of integral
points need not locate any point, whereas a bound on their \emph{height}
immediately gives a complete finite search.  The distinction is relevant before
one asks which induction principles a proof uses.

\subsection{A finite arithmetic interface}

For this discussion, work over $\mathbb Q$ and use a fixed affine presentation.
An input $c$ encodes an integer $m\geq 1$ and finitely many polynomials
$F_1,\ldots,F_r\in\mathbb Z[X_1,\ldots,X_m]$.  It may also contain a finite
certificate $w$ of additional hypotheses.  Signed integers, finite sequences,
and polynomial coefficient arrays are represented by natural-number codes.
Use the concrete signed-integer encoding
\[
 \operatorname{enc}(z)=
 \begin{cases}2z&z\geq0,\\-2z-1&z<0.\end{cases}
\]
Its inverse, absolute value, the elementary ring operations, polynomial
evaluation, sequence length, and sequence entries are primitive recursive.
Choose an injective canonical coding of finite tuples, so each tuple has
exactly one accepted code.  The point predicate rejects noncanonical codes;
equality and list membership therefore compare points rather than multiple
representations of one point.  Malformed codes receive a harmless default
interpretation in the arithmetic operations, so these operations are total
functions on natural numbers.

Let $P(c,p)$ mean that $p$ is a well-formed $m$-tuple of integers and
$F_i(p)=0$ for every $i\leq r$.  Let
\[
 h(c,p)=\max_{1\leq j\leq m}|p_j|
\]
on well-formed tuples.  Both $P$ and $h$ are primitive recursive.  Here $h$ is
an affine coordinate size; no theory of canonical heights is needed.  Finally,
let $G(c)$ be a primitive recursive test for the supplied admissibility
certificate.  We do \emph{not} assume that every desired geometric property has
already been compiled into such a test.  That compilation, its soundness, and
its coverage are separate obligations.

This interface isolates the logical issue without assuming a constructive
theory of algebraic curves.  It also permits statements conditional on a
specific, independently verified certificate.  If a later development chooses
a decidable geometric presentation directly, $G$ can instead be its validity
test.  Conversely, an unbounded, nondecidable hypothesis must not be replaced
by $G$ without justification.

Write
\[
 \operatorname{Bound}(c,B):\equiv
 \forall p\bigl(P(c,p)\longrightarrow h(c,p)\leq B\bigr).
 \tag{\textsf{H}}
\]
The quantifier over $p$ is unbounded.  Consequently,
$\operatorname{Bound}(c,B)$ is generally a $\Pi^0_1$ formula, although every
individual polynomial evaluation and comparison is decidable.  Decidability
of those individual tests does not make the universal assertion decidable.

\subsection{Height bounds and complete lists}

A positive finite-list statement has the form
\[
 \exists L\ \forall p\bigl(P(c,p)\longleftrightarrow p\in L\bigr),
 \tag{\textsf{L}}
\]
where $L$ is a natural-number code for a finite sequence.  Repetitions may be
removed by a primitive recursive filter.  Thus the sequence formulation is
equivalent to a coded duplicate-free list, without invoking arbitrary sets.

There are direct constructive transformations between $(\mathsf L)$ and
$\exists B\,\operatorname{Bound}(c,B)$.  Given $L$, compute the maximum of the
heights of its entries, with maximum $0$ for the empty list.  Every point
occurs in the list, hence its height is at most this maximum.  This direction
only needs the covering implication $P(c,p)\to p\in L$; if the list contains
extraneous tuples, filtering by $P$ removes them.

For the converse, suppose a particular $B$ and a proof of
$\operatorname{Bound}(c,B)$ are available.  Enumerate the box
\[
 \{-B,-B+1,\ldots,B\}^{m}.
\]
It has exactly $(2B+1)^m$ entries.  For example, digits in base $2B+1$
enumerate coordinate tuples, with each digit $d$ interpreted as $d-B$.
Evaluate the defining polynomials on each tuple and retain the tuples for
which every value is zero.  All loops have explicit finite limits.  Their
composition is primitive recursive, including when $m$ and $B$ vary.
Every retained tuple is a point by the filter test.  Every point lies in the
box by the given bound and is therefore retained.  This proves completeness.

The associated iPRA argument uses only elementary coding and induction on
finite loops.  A statement that a loop has inspected its first $k$ entries is
bounded; in the language with all primitive recursive functions it can be
expressed by a quantifier-free characteristic equation.  Thus the usual
induction proof of filter correctness does not require induction over an
arbitrary unbounded formula.  This observation justifies the finite search
\emph{after} a height bound has been proved; it does not supply that bound.

For integral points over a number field, one must additionally specify a
finite encoding of the field and integral elements and give an effective
finite enumeration below the chosen height.  Those additions do not remove
the distinction between bounding the number of points and bounding their
locations.  The rational case already exhibits the logical issue.

\subsection{A bound on the number of points is weaker}

Define $\operatorname{AtMost}(c,N)$ by saying that every list of $N+1$ points
contains a repeated entry.  More explicitly,
\[
 \forall s\Bigl(
   \operatorname{length}(s)=N+1\ \wedge\
   \operatorname{AllPoints}(c,s)
   \ \longrightarrow\
   \operatorname{HasRepeat}(s)\Bigr).
 \tag{\textsf{C}}
\]
The two sequence predicates are primitive recursive: their internal
quantification is bounded by the sequence length.  This is therefore also a
$\Pi^0_1$ assertion.  A complete list gives $(\mathsf C)$ with its length as
$N$, by the finite pigeonhole principle.  A height bound gives it with
$N=(2B+1)^m$.

The reverse implication is not intuitionistically valid in general, even
for primitive recursive predicates.  To see the reason, let $R(e,t)$ mean
that a deterministic program with code $e$, on a fixed input, halts
\emph{for the first time} after exactly $t$ steps.  Bounded simulation makes
$R$ primitive recursive.  Finite computation and determinism prove that
there cannot be two different such times.  In particular the numerical
bound $1$ works uniformly for the number of witnesses of $R(e,\cdot)$.

Nevertheless, no total computable function $b(e)$ can bound every such
halting time.  If it did, simulate program $e$ for $b(e)$ steps.  If it has
halted, answer yes; otherwise the bound guarantees that it never halts, so
answer no.  This would decide the halting problem.  Thus a uniform count
bound, even the constant $1$, does not in general give a computable height
bound or a complete list.

This example concerns arbitrary decidable subsets of the natural numbers.
It is \emph{not} a reduction of the halting problem to integral points on
the curves under discussion.  It supplies no undecidability result for
Siegel's theorem.  Its purpose is narrower: it rules out a purely logical
passage from a cardinality bound to an exhaustive search bound.

Knowing an upper bound $N$ also does not justify searching until $N$ points
have been found.  There may be fewer than $N$ points, or none.  Such a
search would then fail to stop.  A proof providing $N$ actual distinct
points together with an upper bound $N$ is different: the provided points
form a complete list.  The positive witnesses in the lower-bound statement
are essential.

Accordingly, the uniform assertion
\[
 \forall c\bigl(G(c)\to\exists N\,
                         \operatorname{AtMost}(c,N)\bigr)
 \tag{\textsf{CountSiegel}}
\]
would extract a primitive recursive bound on the \emph{number} of points.
It need not solve the effective problem of locating all points.  It would
still require a complete proof in the specified arithmetic.  The existence
of classical quantitative results is not, by itself, such a proof.  The
phrase ``finitely many'' in the requested theorem therefore leaves a real
choice of constructive formulation; this report does not identify
$(\mathsf{CountSiegel})$ with the stronger assertion below.

\subsection{The precise witness-extraction consequence}

For the usual presentation of iPRA, the relevant metatheorem is stronger
than mere recursive realizability.  If
\[
 \mathrm{iPRA}\vdash\forall u\exists v\,A(u,v),
\]
where $A$ is an arbitrary arithmetical formula, then there is a primitive
recursive function symbol $f$ such that
\[
 \mathrm{iPRA}\vdash\forall u\,A(u,f(u)).
 \tag{\textsf{W}}
\]
Salehi records this as Fact~3.5, with a reference to Wehmeier's
primitive-recursive choice theorem
\cite{Salehi2002,Wehmeier1997}.  Multiple input variables are combined by
primitive recursive pairing.  In particular, the theorem is not restricted
to quantifier-free $A$; its applicability to a matrix containing an
unbounded universal quantifier is exactly what is needed here.

Here is the consequence for a hypothetical proof of positive Siegel
finiteness.  Suppose
\[
 \mathrm{iPRA}\vdash
 \forall c\bigl(G(c)\longrightarrow
          \exists B\,\operatorname{Bound}(c,B)\bigr),
 \tag{\textsf{PositiveSiegel}}
\]
and suppose iPRA proves the decidability of $G$.  For any fixed $c$, reason
by that decidability.  If $G(c)$, use $(\mathsf{PositiveSiegel})$ and retain
its witness $B$.  If $\neg G(c)$, take $B=0$; the implication from $G(c)$
is vacuous.  These two cases prove, intuitionistically,
\[
 \mathrm{iPRA}\vdash\forall c\exists B
       \bigl(G(c)\longrightarrow\operatorname{Bound}(c,B)\bigr).
\]
Apply $(\mathsf W)$ with
$A(c,B)\equiv G(c)\to\operatorname{Bound}(c,B)$.  It supplies a primitive
recursive $b$ and an iPRA proof of
\[
 \forall c\bigl(G(c)\longrightarrow
                         \operatorname{Bound}(c,b(c))\bigr).
 \tag{\textsf{ExtractedBound}}
\]
No unrestricted independence-of-premise principle has been used.  The
case distinction on the explicitly decidable certificate predicate is
what permitted moving the existential quantifier.

Composing $b$ with the box enumeration and polynomial filter gives a
primitive recursive procedure which outputs a complete list on every
certified admissible input.  It also decides whether that input has any
integral points, by checking whether the resulting list is empty.  The
numerical bound need not be remotely practical; primitive recursiveness
is a structural conclusion about its construction, not a claim of
feasible running time.

The same argument applies to $(\mathsf{CountSiegel})$, but with
$A(c,N)=G(c)\to\operatorname{AtMost}(c,N)$.  Its extracted function only
supplies $N$.  Substituting this function into a height-bounded search
would be an error: no height bound occurs in the conclusion of the
witness theorem.

There are two further qualifications.  First, when $c$ contains a
certificate, the extracted bound is allowed to depend on that certificate.
Obtaining an algorithm from raw equations requires a procedure that
produces a valid certificate for every admissible raw input.  A
primitive recursive certificate checker alone does not supply such a
procedure.  Second, if admissibility is instead represented by a
nondecidable formula, a realizability argument may require a realizer of
that hypothesis.  It is not legitimate to discard that additional input
and announce a bound depending only on the curve's coefficients.

\subsection{Fixed instances and external schemas}

A proof for one specified curve and a uniform theorem over curve codes
have different strengths.  If iPRA proves
$\exists B\,\operatorname{Bound}(\bar c,B)$ for one fixed numeral
$\bar c$, numerical existence provides a particular standard integer
$b$ with an iPRA proof of $\operatorname{Bound}(\bar c,\bar b)$.
There is no conclusion that the same bound works for other curves.
For that particular curve, however, the proof supplies precisely the
information required for a complete finite search.

There is a subtle intermediate case.  Suppose one asserts externally
that \emph{every} admissible standard curve code $c$ has some iPRA proof
of $\exists B\,\operatorname{Bound}(\bar c,B)$, without presenting a
uniform proof of $(\mathsf{PositiveSiegel})$.  This assertion still has
an effectivity consequence.  On a given admissible input, enumerate iPRA
proofs.  Numerical existence guarantees that eventually one finds a
proof of $\operatorname{Bound}(\bar c,\bar b)$ for some numeral
$\bar b$.  Return $b$ and perform the finite search.  Proof enumeration
and recognition are effective, so this is a partial recursive algorithm
which terminates on every admissible input covered by the asserted
schema.  Its running time, and the resulting bound as a function of $c$,
need not be primitive recursive.

Thus a nonuniform assertion of \emph{provability of all instances}
must not be confused with bare classical existence of a bound for
each instance.  The latter does not furnish recognizable iPRA proofs
to enumerate.  Nor does a collection of proofs for selected examples
establish that all admissible instances are covered.  These distinctions
prevent both an exaggerated impossibility claim and an understated
claim about what a proposed constructive theorem would accomplish.

\subsection{Negations, translations, and hidden induction}

A possible weaker target is
\[
 \forall c\bigl(G(c)\longrightarrow
              \neg\neg\exists B\,\operatorname{Bound}(c,B)\bigr).
 \tag{\textsf{NegativeSiegel}}
\]
The witness theorem does not remove the double negation.  In particular,
ordinary Markov-style reasoning for a decidable search predicate does
not apply directly: the predicate being searched for is
$\operatorname{Bound}(c,B)$, an unbounded universal assertion.  Checking
a large finite region cannot verify it.

The statement that points are not unbounded,
\[
 \neg\forall B\exists p\,
                    \bigl(P(c,p)\wedge h(c,p)>B\bigr),
\]
is another possible negative formulation.  It follows from positive
boundedness and also from its double negation.  One should not assume
that all these negative formulations are intuitionistically equivalent;
interchanging negation and unbounded quantifiers itself requires
justification.  Likewise, the absence of a specified infinite sequence
of distinct points is not a coded complete list of the points.

Negative translation is useful only with the correct source theory.
A classical \emph{PRA} proof can be translated into an iPRA proof of
its negative translation, because the relevant defining equations and
quantifier-free induction are preserved.  A proof in arbitrary classical
mathematics does not automatically become an iPRA proof by the same
operation.  Its induction, choice, comprehension, and algebraic
existence principles must first be accounted for.  Therefore a classical
proof of Siegel's theorem cannot simply be relabeled a proof of
$(\mathsf{NegativeSiegel})$ in iPRA.

Similar care is needed with formalization.  A Lean proof about
$\mathbb Z$, polynomials, and finite lists can establish the correctness
of a finite search.  A Lean theorem whose hypothesis already asserts a
height bound is a conditional result.  Neither one, by itself, proves
that iPRA derives the required height bound.  For the latter claim,
one needs an explicit translation to object-language arithmetic and a
derivation whose induction formulas satisfy the stated restriction.
Using Lean's unrestricted natural-number induction in the metatheory
does not discharge that obligation.

\subsection{A finite kernel that can be completed independently}

Several useful lemmas do not depend on the unresolved global argument.
First, primitive recursive evaluation of a polynomial code agrees with
the corresponding finite sum of monomials.  A proof proceeds through
the coefficient list and then through each exponent vector; all
intermediate identities are ring equations.  Second, decoding the
base-$(2B+1)$ expansion of an index below $(2B+1)^m$ gives exactly one
tuple in the box.  Conversely, encoding the shifted coordinates of
any box tuple returns such an index.  Both assertions reduce to
bounded division and remainder identities.

Third, filtering a finite sequence by the point predicate is sound,
and every original entry satisfying that predicate occurs in the
filtered sequence.  Fourth, the maximum-height operation bounds the
height of each list entry.  Finally, the pigeonhole argument converts
a complete list into a bound on the number of distinct points.  These
proofs may be organized with characteristic functions for bounded
sequence predicates, so their loop induction remains quantifier-free.
Their correctness is uniform in the coefficient data, the number of
coordinates, and the supplied numerical bound.

These results constitute a genuine computational kernel, but its
interface must retain the hypothesis
$\operatorname{Bound}(c,B)$.  An executable routine that receives $B$
and returns the filtered box list is already useful; calling its
output ``all integral points'' without verifying that hypothesis is
not justified.  Likewise, checking finitely many examples of the
routine checks implementation behavior, not the global theorem.

An arithmetic proof checker can instead check a finite derivation of
the universal bound.  This is a stronger certificate than merely
testing points up to a chosen cutoff.  Its syntax, inference rules,
and induction restriction must be specified and its soundness
established in the metatheory.  Having such a checker does not create
the required derivation.  Separating the finite search kernel, the
proof checker, and the mathematical height-bound argument prevents a
correct piece of infrastructure from concealing the missing theorem.

\subsection{What this audit establishes}

The established logical conclusion is conditional and precise: a
uniform iPRA proof of positive, height-bounded Siegel finiteness on a
decidably certified family yields a primitive recursive height bound
and complete enumeration on that family.  A uniform iPRA proof of
cardinality finiteness yields only a primitive recursive cardinality
bound.  An admissible negative formulation has still different
computational content.

None of these observations proves that Siegel's theorem, in any of
these formulations, is unprovable in iPRA.  They identify what must be
proved and what cannot be supplied by a purely formal change of
language.  In particular, stating that the positive version would
settle an unresolved effective problem is not an independence theorem.
A completed contribution would have to provide the missing arithmetic
argument, or explicitly choose and prove a weaker constructive
formulation.  The finite coding and extraction results above remain
useful in either direction, while making the remaining mathematical
work visible.
